\section*{Fachliche Kurzfassung}
\addcontentsline{toc}{section}{Fachliche Kurzfassung}

Die vorliegende Arbeit evaluiert systematisch mehrere Deep-Learning-Architekturen für die automatische Klassifikation und Detektion von Recyclingmaterialien auf ressourcenbeschränkter Edge-Hardware \autocite{yosinski2014transferable}.
Ein Vergleich eines Scratch-CNNs, eines vortrainierten MobileNetV2 mit Transfer Learning sowie eines YOLO11n-Modells zeigt, dass Post-Training-Quantisierung (INT8) die Modellgröße des fine-tuned MobileNetV2 um einen Faktor 9 auf \SI{2.61}{\mega\byte} komprimiert, bei einem vernachlässigbaren Genauigkeitsverlust von 0,07 Prozentpunkten \autocite{krishnamoorthi2018quantizing,deng2020model}.
Das fine-tuned MobileNetV2 erreicht eine Validierungsgenauigkeit von \SI{87.42}{\percent} und übertraf das Scratch-CNN um \num{26.42} Prozentpunkte; YOLO11n erzielt auf der besten Datensatzkombination \SI{60.7}{\percent} mAP50 \autocite{ultralyticsyolov8,proenca2020taco,yang2016classification}.
Die Ergebnisse deuten auf die Machbarkeit echtzeitfähiger, kostengünstiger Sortiersysteme hin, die auf ressourcenbeschränkter Edge-Hardware (Raspberry Pi Zero 2W oder Raspberry Pi 4) betrieben werden können.
